% Candidate: different calculation and first-wild character recovery.
\subsection{The cyclotomic different closes the remaining first-wild comparison}

The following field lemma is independent of any formal group.
The residue field is allowed to be algebraically closed and infinite.
All different exponents and lower valuations are normalized on the
top field named in their subscript. Upper breaks use the conventions
of B60 and B97, with no conductor-plus-one shift.

\begin{lemma}[A lower bound for the second cyclotomic upper break]
\label{lem:B-RD98-different}
\status{Candidate; complete proof and source checks awaiting audits.}
Let $p\ge5$ and let $K$ be complete discretely valued of mixed
characteristic $(0,p)$ with algebraically closed residue field,
$e_K=p(p-1)$ and $\mu_p\subset K$, but $\mu_{p^2}\not\subset K$.
Then $C=K(\mu_{p^3})/K$ is cyclic of degree $p^2$. If its two
upper breaks are $\delta_1<\delta_2$, then
\[
            \delta_1+p\delta_2\ge p(p^2-p-1).
                                                        \tag{B-RD98.1}
\]
Consequently, if $\delta_1\le p-1$, then
\[
                     \delta_2\ge p^2-p-2+\frac1p.
                                                        \tag{B-RD98.2}
\]
No integrality theorem for upper breaks is needed.
\end{lemma}
\begin{proof}
Choose a Cohen coefficient ring $\Lambda\subset\mathcal O_K$
and put $K_0=\operatorname{Frac}(\Lambda)$, as in the full
proof of Lemma~\ref{lem:B-RD86-cyclotomic}. Its residue field
is the same perfect field and its uniformizer is $p$.
For completeness, $[K:K_0]=e_K$: if $\pi$ is a uniformizer
of $K$, the elements $1,\pi,\ldots,\pi^{e_K-1}$ give a basis
of $\mathcal O_K/p\mathcal O_K$ over the residue field.
Lifting coefficients to $\Lambda$ and iterating $p$-adically
shows that they span $\mathcal O_K$ over $\Lambda$.
Their distinct valuations modulo $e_K$ prove independence.
This also proves the required finiteness over the chosen Cohen
field without a finite-residue assumption.

Set
\[
 K_1=K_0(\mu_p),\qquad C_0=K_0(\mu_{p^3}),\qquad C=KC_0.
                                                        \tag{B-RD98.3}
\]
Cyclotomic Eisenstein polynomials give
$[K_1:K_0]=p-1$, $[K:K_1]=p$, and $C_0/K_1$ cyclic of
degree $p^2$. Its unique degree-$p$ subfield is
$K_0(\mu_{p^2})$. Therefore $K\cap C_0=K_1$: a larger
intersection would equal the degree-$p$ field $K$ and force
$\mu_{p^2}\subset K$, contrary to hypothesis. It follows that
\[
                       [C:K]=p^2,\qquad[C:C_0]=p.
                                                        \tag{B-RD98.4}
\]
All these extensions have unchanged residue field and are
totally ramified. The first assertion of the lemma follows.

Write $d(A/B)=v_A(\mathfrak D_{A/B})$ for a different exponent.
The two elementary estimates needed are
\[
                      d(K/K_1)\le p^2-1,
                      \qquad d(C/C_0)\ge p-1.
                                                        \tag{B-RD98.5}
\]
To prove the first, a uniformizer $\pi$ of $K$ generates its
ring over $\mathcal O_{K_1}$ by the same uniformizer-expansion
argument. Let its Eisenstein polynomial be
$f(X)=X^p+\sum_{i=0}^{p-1}a_iX^i$.
The different is generated by $f'(\pi)$, by the full
monogenic different proof in \cite[Proposition3.18 and proof,
pp.~17--18]{ColmezCorpsLocaux}. The leading derivative term
has valuation
\[
                    v_K(p\pi^{p-1})=p^2-1.
                                                        \tag{B-RD98.6}
\]
Every other nonzero derivative term $i a_i\pi^{i-1}$,
$1\le i\le p-1$, has valuation congruent to $i-1$ modulo $p$,
because $v_K(a_i)$ is divisible by $p$ and $i$ is a unit.
These residues are pairwise distinct and differ from $-1$,
the residue of (B-RD98.6). Thus no cancellation can remove the
least term, and the first estimate follows. For $C/C_0$ the
same monogenic Eisenstein derivative formula shows that every
term has valuation at least $p-1$, proving the second estimate.

The known cyclotomic extension $C_0/K_1$ has lower breaks
$p-1$ and $p^2-1$. Indeed $\zeta_{p^3}-1$ is a uniformizer;
an element whose cyclotomic exponent is $1+pa$ with $p\nmid a$
has uniformizer displacement of valuation $p$, whereas a nonzero
exponent $1+p^2a$ gives valuation $p^2$.
The different formula, obtained by summing these displacement
valuations, therefore gives
\[
                         d(C_0/K_1)=2p^2(p-1).
                                                        \tag{B-RD98.7}
\]
This is also the integral ramification formula in
\cite[Corollary3.19 and proof, p.~18]{ColmezCorpsLocaux}.
Different transitivity, whose complete trace-dual proof is
\cite[Proposition3.17 and proof, p.~17]{ColmezCorpsLocaux},
applied to the two actual towers over $K_1$ gives
\[
\begin{split}
 d(C/K)+p^2d(K/K_1)
   &=d(C/C_0)+p\,d(C_0/K_1),\\
 d(C/K)&\ge p-1+2p^3(p-1)-p^2(p^2-1)\\
       &=p^2(p-1)^2+p-1.
\end{split}                                             \tag{B-RD98.8}
\]
The ramification indices, not just the field names, determine
the factors $p^2$ and $p$ in this identity.

A cyclic group of order $p^2$ has two positive lower breaks:
the $p$th-power subgroup has larger break because its first
ramification coefficient vanishes on taking a $p$th power.
They are $\delta_1$ and $p\delta_2-(p-1)\delta_1$.
Summing its lower ramification groups, or equivalently applying
the same different formula, yields
\[
              d(C/K)=(p-1)(\delta_1+p\delta_2+p+1).
                                                        \tag{B-RD98.9}
\]
Substitute (B-RD98.8), divide by $p-1$, and subtract $p+1$.
This proves (B-RD98.1). Using $\delta_1\le p-1$ gives
(B-RD98.2), with its rational endpoint retained. No Hasse--Arf
integrality or finite-residue class-field theorem is invoked.
\end{proof}

\begin{theorem}[Irreducible residual recovery at the first wild index]
\label{thm:B-RD98-irreducible}
\status{Candidate; complete proof and source checks awaiting audits.}
For every $p\ge5$, over a field satisfying the hypotheses
$e_K=p(p-1)$ and $\mu_p\subset K$, an actual $p^3$-torsion
isomorphism between two whole potentially supersingular
height-two groups with irreducible residual inertia determines
their finite special characters.
\end{theorem}
\begin{proof}
Suppose the characters differ and apply the complete B97
actual point-field results. Put $q=p^2$, $N=q-1$, let $L/K$
be the common tame good field of degree $N$, and let $F_2,F_3$
be the actual common division fields. Then
\[
\begin{gathered}
 r_1=\dim W_1\in\{3,4\},\qquad
       \operatorname{tr}|_{W_1}\ne0,\\
 c_1=N(p-1)=\text{last upper break of }F_2/L,\qquad
 c_2=(q^2-1)(p-1),\qquad
 \text{last upper break of }F_3/F_2\le p^{r_1-2}c_2.
\end{gathered}                                          \tag{B-RD98.10}
\]
The last assertion includes the full terminal character space:
if $P_1^+,P_2^+$ form an actual $p^3$-basis and
$P_i=pP_i^+$, then $E_i=L(P_i)$, $T_i=L(P_i^+)$ satisfy
$T_i/E_i$ Galois of degree $q$, $F_2\cap T_i=E_i$ and
$F_2T_1T_2=F_3$. On the terminal subgroup, the two actual
translation maps are $B\mapsto Be_1$ and $B\mapsto Be_2$.
Their kernels have zero intersection. Thus the two corresponding
degree-$q$ quotient character spaces span every terminal
character, including the actual determinant. Their conductor
bounds come from $\psi_{F_2/E_i}(c_2)\le p^{r_1-2}c_2$,
not from an identification of unrelated cyclic extensions.

On $L$ the determinant is the cyclotomic character. Its
nonzero first-layer trace means $\mu_{p^2}\not\subset L$,
and therefore $\mu_{p^2}\not\subset K$. The field lemma
gives $C=K(\mu_{p^3})/K$ cyclic of degree $p^2$.
Its first degree-$p$ subfield is $K(\mu_{p^2})$.
Tame base change of degree $N$ multiplies every lower and
upper break by $N$: the uniformizer/Hensel argument of B86
scales each lower displacement by $N$, and the Herbrand
integrals scale their arguments by the same factor.
Thus the two upper breaks of $LC/L$ are $N\delta_1,N\delta_2$.

The actual first cyclotomic layer $L(\mu_{p^2})$ is a quotient
of $F_2/L$. Consequently
\[
                         N\delta_1\le c_1,
                   \qquad \delta_1\le p-1.
                                                        \tag{B-RD98.11}
\]
By (B-RD98.2) its next break $d_2=N\delta_2$ satisfies
\[
 d_2\ge N\left(p^2-p-2+\frac1p\right)>c_1,
 \qquad q(d_2-c_1)-c_2
          \ge N\bigl(p^4-3p^3+1\bigr)>0.
                                                        \tag{B-RD98.12}
\]
Both strict inequalities hold for every $p\ge5$.

We identify the restriction field precisely. The actual
$LC/L$ is cyclic of degree $p^2$, whereas $F_2/L$ is
elementary abelian and contains its first cyclotomic layer.
Hence their intersection is exactly $L(\mu_{p^2})$.
The extension $F_2C/F_2$ is therefore cyclic of degree $p$
and is a nontrivial quotient of the actual terminal field
$F_3/F_2$. Because every upper break of $F_2/L$ is at most
$c_1<d_2$, upper quotient compatibility and subgroup
conversion give its conductor exactly
\[
 \psi_{F_2/L}(d_2)
  =\psi_{F_2/L}(c_1)+p^{r_1}(d_2-c_1)
  >p^{r_1}(d_2-c_1).
                                                        \tag{B-RD98.13}
\]
For an explicit filtration check, above $c_1$ the ramification
group of the actual compositum $F_2C/L$ lies in its subgroup
over $F_2$; its image in $LC/L$ is the nontrivial last cyclic
line until $d_2$. Thus (B-RD98.13) concerns the actual restricted
cyclotomic character, including its nontrivial common first layer.

The final inequality of (B-RD98.12) makes (B-RD98.13) strictly
larger than $p^{r_1-2}c_2$, contrary to (B-RD98.10).
This contradiction eliminates the remaining traceful case.
The traceless case and the smaller first layers were already
excluded by B97, so the theorem follows.
\end{proof}

\begin{corollary}[The exact first wild height-two index at five]
\label{cor:B-RD98-exact-five}
\status{Candidate; complete proof and source checks awaiting audits.}
In the full B.4 height-two setting,
\[
                         \nu_{\min}(5,20,2)=2.
                                                        \tag{B-RD98.14}
\]
Thus actual $125$-torsion determines the requested prime-to-five
special characters for every allowed group and base of this index,
whereas actual $25$-torsion does not do so uniformly.
\end{corollary}
\begin{proof}
The full determinant/dimension and whole-versus-filtered
reductions of B76 apply at $p=5,e=20,n=3$: the determinant
condition is $n\ge v_5(e)+2$, and the mixed exclusion condition
is $20\le5^{2n-3-v_5(e)}=25$. Hence any failure at actual
$125$-torsion must compare two whole potentially supersingular
height-two groups. Passing to the completed maximal unramified
extension preserves the absolute index and the inertia character
question, so one may use the algebraically closed residue field
in the preceding results.

If their residual representation is irreducible, the theorem
applies. If it is reducible with nontrivial residual wild action,
the full actual-field exclusion of B96 applies. The remaining
tame reducible residual case, including the earlier scalar and
trivial subcases, was excluded by B82--B84 and B87--B90.
These exhaust the actual height-two possibilities and prove
determination at $125$.

For the strict lower bound, Theorem~\ref{thm:B-RD63-arc}
applies at $p=5$, actual level $p^2=25$, and $e=20>5$.
It constructs an allowed field and a whole-versus-filtered
comparison with different requested finite special characters.
Thus index one fails. The manuscript index is $\nu=n-1$
for actual level $p^n$; the upper and lower bounds therefore
give exactly (B-RD98.14).
\end{proof}

\begin{remark}[Scope, source proof and intended exact supplement]
\label{rem:B-RD98-scope}
\status{Candidate; no least-index promotion before all audits.}
The all-prime theorem concerns the irreducible-residual
whole-group branch at $e=p(p-1)$. The exact full B.4 conclusion
is asserted here only for $(p,e,h)=(5,20,2)$, using the proved
five-specific reducible branches. No exact first-wild index for
another prime, higher height or a general ramification index
is inferred.

For this proof the retained original Colmez text was freshly
read through the complete Proposition3.17 trace-dual transitivity
proof, Proposition3.18 monogenic derivative proof and its
continuation, and Corollary3.19 ramification-different proof,
printed pp.~17--18. The arbitrary-perfect-residue field scope
and the complete B86 Cohen-field argument were also checked.
No new external web opening or image inspection is claimed.
The complete actual B97 point-field, normality, intersection
and terminal-character spanning proofs are an explicit dependency.
A separate exact program will check the symbolic different and
conductor identities without replacing any field-theoretic step.
\end{remark}
